\documentclass[12pt,a4paper]{article}
\usepackage[utf8]{inputenc}
\usepackage[T1]{fontenc}
\usepackage[french]{babel}
\usepackage{geometry}
\usepackage{xcolor}
\usepackage{tcolorbox}
\usepackage{fancyhdr}

% --- DESIGN & COULEURS ---

% Police sans serif (plus propre pour une fiche)
\renewcommand{\familydefault}{\sfdefault}

% Palette de couleurs
\definecolor{mainBlue}{RGB}{0, 70, 100}      % Bleu foncé (Titres principaux)
\definecolor{softTeal}{RGB}{0, 140, 160}     % Cyan (Sous-titres / Définitions)
\definecolor{alertRed}{RGB}{180, 40, 40}     % Rouge (Attention / Important)
\definecolor{bgGray}{RGB}{245, 247, 250}     % Gris très clair (Fond des boîtes)

% Configuration des boites (tcolorbox)
\tcbset{
    colback=bgGray,           % Couleur de fond
    colframe=mainBlue,        % Couleur du cadre
    coltitle=white,           % Couleur du titre
    fonttitle=\bfseries\large,% Police du titre
    arc=3mm,                  % Arrondi des coins
    boxrule=0.6mm,            % Épaisseur du trait
    left=2mm, right=2mm, top=2mm, bottom=2mm, % Marges internes
    after skip=0.4cm          % Espace après la boite
}

% Commandes rapides pour le texte
\newcommand{\notion}[1]{\textbf{\textcolor{mainBlue}{#1}}}
\newcommand{\important}[1]{\textit{\textcolor{alertRed}{#1}}}
\newcommand{\definition}[1]{\textbf{\textcolor{softTeal}{#1}}}

\begin{document}





\section*{Problématique}
\textbf{Comment s'établissent les relations entre l'entreprise et son environnement économique ?}

\section{Définitions essentielles}

\begin{tcolorbox}[title=\textbf{DÉFINITION : Agent économique}]
Individu ou ensemble d'individus constituant un centre de décision économique indépendant, regroupés selon leur fonction principale.
\end{tcolorbox}

\vspace{3mm}
\begin{tcolorbox}[title=\textbf{DÉFINITION : Production marchande}]
Biens et services vendus sur un marché à un prix concurrentiel, assurée par les entreprises.
\end{tcolorbox}

\vspace{3mm}
\begin{tcolorbox}[title=\textbf{DÉFINITION : Production non marchande}]
Biens et services fournis gratuitement ou à des prix non significatifs (< 50\% du coût), assurée par les administrations publiques.
\end{tcolorbox}

\vspace{3mm}
\begin{tcolorbox}[title=\textbf{DÉFINITION : Circuit économique}]
Modèle schématique des échanges et interdépendances entre agents économiques, représentant flux réels et monétaires.
\end{tcolorbox}

\vspace{3mm}
\begin{tcolorbox}[title=\textbf{DÉFINITION : Marché des capitaux}]
Lieu de rencontre entre offreurs (épargnants) et demandeurs de capitaux (entreprises), où le prix du capital est le taux d'intérêt.
\end{tcolorbox}

\section{Les 4 agents économiques principaux}

\subsection{Ménages}
\begin{tcolorbox}[title=\textbf{POINT CLÉ}]
Fonction principale : CONSOMMATION de biens et services pour satisfaire directement leurs besoins
\end{tcolorbox}
\begin{itemize}
    \item Individus ou groupes vivant dans un même logement
    \item Offrent leur travail sur le marché du travail
    \item Consomment et épargnent
\end{itemize}

\subsection{Entreprises}
\begin{tcolorbox}[title=\textbf{POINT CLÉ}]
Fonction principale : PRODUCTION de biens et/ou services pour les vendre et réaliser un profit
\end{tcolorbox}
\begin{itemize}
    \item Production marchande destinée à la vente
    \item Demandeurs de travail et de capitaux
    \item Relations B2B avec d'autres entreprises
\end{itemize}

\subsection{Banques}
\begin{tcolorbox}[title=\textbf{POINT CLÉ}]
Fonction principale : FINANCEMENT de l'économie par collecte et octroi de crédits
\end{tcolorbox}
\begin{itemize}
    \item Intermédiaires financiers
    \item Collectent l'épargne des ménages
    \item Accordent des crédits aux entreprises et ménages
\end{itemize}

\subsection{Administrations publiques}
\begin{tcolorbox}[title=\textbf{POINT CLÉ}]
Fonction principale : Répondre aux besoins d'intérêt général et produire des services non marchands
\end{tcolorbox}
\begin{itemize}
    \item État, Sécurité sociale, collectivités locales
    \item Fournissent services publics (éducation, sécurité, justice)
    \item Financés par impôts, taxes, cotisations sociales
\end{itemize}

\section{Les principaux marchés}

\subsection{Marché des biens et services}
\begin{tcolorbox}[title=\textbf{POINT CLÉ}]
Échange : Biens/services (flux réels) contre Paiements (flux monétaires)
\end{tcolorbox}
\begin{itemize}
    \item Entreprises → Ménages : produits et services
    \item Ménages → Entreprises : chiffre d'affaires
    \item Décisions influencées par la conjoncture économique
\end{itemize}

\subsection{Marché du travail}
\begin{tcolorbox}[title=\textbf{POINT CLÉ}]
Échange : Travail (flux réels) contre Salaires (flux monétaires)
\end{tcolorbox}
\begin{itemize}
    \item Ménages → Entreprises : force de travail
    \item Entreprises → Ménages : rémunération
    \item Rencontre offre et demande de travail
\end{itemize}

\subsection{Marché financier}
\begin{tcolorbox}[title=\textbf{POINT CLÉ}]
Lieu d'émission et d'échange des valeurs mobilières (actions, obligations)
\end{tcolorbox}
\begin{itemize}
    \item \textbf{Marché primaire} : émission de nouvelles valeurs
    \item \textbf{Marché secondaire} : échange de titres existants
    \item Financement externe direct (sans intermédiaire bancaire)
\end{itemize}

\section{Relations entreprise-État}

\begin{tcolorbox}[title=\textbf{POINT CLÉ}]
L'État fournit services non marchands, prestations sociales et subventions
\end{tcolorbox}

\subsection{Ce que l'État apporte aux entreprises}
\begin{itemize}
    \item Services publics (infrastructures, éducation, sécurité)
    \item Subventions et aides publiques
    \item Cadre juridique et réglementaire
\end{itemize}

\begin{tcolorbox}[title=\textbf{EXEMPLE}]
Bpifrance : banque publique qui finance et accompagne les projets d'entreprises (innovation, international, fonds propres)
\end{tcolorbox}

\subsection{Ce que les entreprises apportent à l'État}
\begin{itemize}
    \item Impôt sur les sociétés
    \item Cotisations sociales
    \item Taxes diverses
    \item Création d'emplois et de richesse
\end{itemize}

\section{Points clés à retenir}

\begin{tcolorbox}[title=\textbf{POINT CLÉ}]
L'entreprise est au centre d'un réseau de relations avec tous les agents économiques
\end{tcolorbox}

\begin{tcolorbox}[title=\textbf{POINT CLÉ}]
Ces relations créent des flux réels (biens, services, travail) et des flux monétaires (paiements, salaires)
\end{tcolorbox}

\begin{tcolorbox}[title=\textbf{POINT CLÉ}]
La conjoncture économique influence les comportements de consommation et d'investissement
\end{tcolorbox}

\begin{tcolorbox}[title=\textbf{POINT CLÉ}]
Le numérique transforme les relations client et nécessite de nouvelles approches marketing
\end{tcolorbox}

\begin{tcolorbox}[title=\textbf{POINT CLÉ}]
Les entreprises collaborent avec leurs fournisseurs pour améliorer performance et création de valeur
\end{tcolorbox}

\section{Mots-clés à connaître}
\begin{itemize}
    \item Agents économiques
    \item Production marchande/non marchande  
    \item Circuit économique
    \item Flux réels/monétaires
    \item Marché des biens et services
    \item Marché du travail
    \item Marché des capitaux
    \item Marché financier
    \item Conjoncture économique
    \item Financement externe direct/indirect
\end{itemize}

\end{document}